%!TEX TS-program = xelatex
\documentclass[12pt,a4paper]{article}
\usepackage{fontspec}
\setmainfont{Times New Roman}
\setmonofont{Consolas}
\usepackage[russian]{babel}
\usepackage[left=25mm,right=15mm,top=20mm,bottom=20mm]{geometry}
\usepackage{amsmath,array,longtable,listings}

\pagestyle{plain}

\setlength{\parindent}{1.25cm}
\setlength{\parskip}{3pt}
\lstset{basicstyle=\ttfamily\fontsize{9}{10.5}\selectfont,
 columns=fullflexible,keepspaces=true,breaklines=true,
 numbers=none,tabsize=4,
 showstringspaces=false,frame=single,xleftmargin=12pt}
\begin{document}
\begin{center}

\end{center}
\section*{1. Цель работы}
Изучить основные команды ассемблера, работу с регистрами, памятью,
символьными данными и консольными функциями ввода-вывода Windows.
\section*{2. Постановка задачи}
Ввести целое число $x$ в восьмеричной системе и вывести его в десятичной.
Вычислить полином
\[
 y=11x^2-2x+1
\]
и вывести $y$ в двоичной и десятичной системах. Промежуточные данные
сохранять в памяти. Не использовать \texttt{wsprintfA}; ввод и вывод
сопровождать поясняющими сообщениями.
	
Программа написана на MASM для Win32, модель памяти FLAT, соглашение STDCALL.
Ввод допускает цифры 0--7 и необязательный начальный минус.
По условию вычисления не должны вызывать переполнение; используемый
диапазон $-13972\leq x\leq13972$ обеспечивает это для всех промежуточных
значений. Некорректный ввод и переполнение выявляются с сообщением об ошибке.
\section*{3. Алгоритм}
\begin{enumerate}
\item Получить дескрипторы консоли через \texttt{GetStdHandle}.
Вывести приглашение и прочитать строку функцией \texttt{ReadConsoleA}.
Исключить завершающие CR и LF, учесть знак.
\item Преобразовать восьмеричные цифры слева направо по формуле
$v\leftarrow 8v+d$, начиная с $v=0$; сохранить каждое новое значение в памяти.
При необходимости изменить знак и записать число в \texttt{x}.
\item Вывести $x$ в десятичной системе. Для преобразования делить модуль
числа на 10, сохранять частное и остаток в памяти, записывать цифры в буфер
справа налево. Для отрицательного числа добавить минус.
\item Вычислить и сохранить $x^2$, $11x^2$, $-2x$, $11x^2-2x$ и $y$ в
\texttt{square}, \texttt{quadratic}, \texttt{linear}, \texttt{partial},
\texttt{result} соответственно.
\item Аналогично шагу 3 преобразовать $y$ сначала с основанием 2, затем 10.
Вывести пояснения и строки через \texttt{WriteConsoleA}.
Завершить программу функцией \texttt{ExitProcess}.
\end{enumerate}
Цикл вывода выполняется хотя бы один раз, поэтому ноль тоже преобразуется
корректно. Буфер заполняется с конца: отдельное обращение строки не требуется.
\clearpage
\section*{4. Результаты тестирования}
Проверка выполнена запуском \texttt{lab1.exe} в консоли Windows.
Входные строки задаются в восьмеричной системе: например,
\texttt{10} означает $10_8=8_{10}$. Ожидаемые значения вычислены
по формуле $y=11x^2-2x+1$ и сопоставлены с сохранёнными результатами запусков.

\begin{center}
\textbf{Таблица 1 — Проверка вычисления полинома}
\small
\setlength{\tabcolsep}{4pt}
\renewcommand{\arraystretch}{1.3}
\begin{tabular}{|c|r|r|r|r|p{37mm}|}
\hline
№ & \shortstack{Ввод\\(основание 8)} & \shortstack{Вывод\\$x_{10}$} & \shortstack{Ожидаемое\\$y_{10}$} & \shortstack{Полученное\\$y_{10}$} & Назначение теста\\\hline
1 & \texttt{0} & 0 & 1 & 1 & Нулевое значение \\\hline
2 & \texttt{1} & 1 & 10 & 10 & Положительная единица \\\hline
3 & \texttt{-1} & -1 & 14 & 14 & Отрицательная единица \\\hline
4 & \texttt{7} & 7 & 526 & 526 & Цифра 7 \\\hline
5 & \texttt{10} & 8 & 689 & 689 & Переход к двум разрядам \\\hline
6 & \texttt{-10} & -8 & 721 & 721 & Отрицательное число \\\hline
7 & \texttt{123} & 83 & 75614 & 75614 & Многозначное число \\\hline
8 & \texttt{00017} & 15 & 2446 & 2446 & Ведущие нули \\\hline
9 & \texttt{33224} & 13972 & 2147356681 & 2147356681 & Верхняя граница \\\hline
10 & \texttt{-33224} & -13972 & 2147412569 & 2147412569 & Нижняя граница \\\hline
\end{tabular}
\end{center}

\begin{center}
\textbf{Таблица 2 — Проверка двоичного вывода результата}
\small
\setlength{\tabcolsep}{6pt}
\renewcommand{\arraystretch}{1.2}
\begin{tabular}{|c|l|l|}
\hline
№ теста & Ожидаемое $y_2$ & Полученное $y_2$\\\hline
1 & \texttt{1} & \texttt{1} \\\hline
2 & \texttt{1010} & \texttt{1010} \\\hline
3 & \texttt{1110} & \texttt{1110} \\\hline
4 & \texttt{1000001110} & \texttt{1000001110} \\\hline
5 & \texttt{1010110001} & \texttt{1010110001} \\\hline
6 & \texttt{1011010001} & \texttt{1011010001} \\\hline
7 & \texttt{10010011101011110} & \texttt{10010011101011110} \\\hline
8 & \texttt{100110001110} & \texttt{100110001110} \\\hline
9 & \texttt{1111111111111100001000000001001} & \texttt{1111111111111100001000000001001} \\\hline
10 & \texttt{1111111111111101110101001011001} & \texttt{1111111111111101110101001011001} \\\hline
\end{tabular}
\end{center}
Во всех десяти тестах десятичный и двоичный результаты совпали
с ожидаемыми. Код завершения программы — 0.

\clearpage
\begin{center}
\textbf{Таблица 3 — Проверка обработки ошибок}
\small
\setlength{\tabcolsep}{5pt}
\renewcommand{\arraystretch}{1.3}
\begin{tabular}{|c|l|p{40mm}|p{35mm}|p{35mm}|}
\hline
№ & Ввод & Назначение теста & Ожидаемая реакция & Полученная реакция\\\hline
11 & \texttt{8} & Недопустимая цифра & Ошибка, код 1 & Ошибка, код 1\\\hline
12 & \texttt{-} & Отсутствие цифр & Ошибка, код 1 & Ошибка, код 1\\\hline
13 & \texttt{12a} & Недопустимая буква & Ошибка, код 1 & Ошибка, код 1\\\hline
14 & \texttt{33225} & Переполнение при вычислении полинома & Ошибка, код 1 & Ошибка, код 1\\\hline
\end{tabular}
\end{center}
Во всех четырёх случаях выведено сообщение
\texttt{Error: invalid octal input or 32-bit overflow.}
и получен ожидаемый код завершения 1.

Пример корректного запуска при $x=10_8=8_{10}$:
\begin{lstlisting}[numbers=none]
Enter x in base 8 (digits 0-7, optional minus): 10
x (decimal) = 8
11*x*x - 2*x + 1 (binary) = 1010110001
11*x*x - 2*x + 1 (decimal) = 689
\end{lstlisting}
Все 14 тестов пройдены.

\section*{5. Текст программы}
\lstinputlisting{../code/main.asm}
\section*{6. Вывод}
Реализованы ввод восьмеричного числа, вычисление полинома и ручное
преобразование чисел для вывода. Проверены ноль, оба знака входного числа,
ведущие нули, границы диапазона и обработка ошибок. Результаты тестов
подтверждают правильность работы программы.
\end{document}

